% !TEX root = Tomato Soup.tex
\input{../preamble.tex}

\begin{document}

% ================= PAGE 1 =================
\makeRecipeTitle{Tomato Soup}

\section*{Ingredients}

\begin{itemize}
      \item Unsalted Butter – 2 Tbsp
      \item Yellow Onion, finely chopped – 1 large (about 3 cups)
      \item Garlic, minced – 1.5 cloves (1 Tbsp)
      \item Crushed Tomatoes (with juice, preferably San Marzano) – two 28‑ounce cans
      \item Chicken Stock – 1 cup
      \item Chopped Fresh Basil – 2 Tbsp (about 1/8 cup)
      \item Sugar – 1/2 Tbsp, or to taste
      \item Black Pepper – 1/4 tsp, or to taste
      \item Heavy Whipping Cream – 1/4 cup
      \item Parmesan Cheese, grated – 3 Tbsp, or to taste
\end{itemize}

\newpage
\makeRecipeTitle{Tomato Soup}
\section*{Instructions}

\begin{enumerate}

      \item \textbf{Sauté Onion and Garlic} \\
            Melt the butter in a nonreactive pot or enameled Dutch oven over medium heat. Add the chopped onions and sauté for 10–20 minutes, stirring occasionally, until softened and golden. Add the minced garlic and sauté for 1 minute until fragrant.

      \item \textbf{Simmer the Soup} \\
            Stir in the crushed tomatoes with their juice, chicken stock, chopped basil, sugar, and black pepper. Bring the soup to a boil, then immediately reduce the heat to low. Partially cover with a lid and let it simmer for 10 minutes.

      \item \textbf{Blend to Desired Consistency} \\
            Remove the pot from the heat. If you prefer a chunky soup, leave it as is. For a creamier texture, use an immersion blender directly in the pot. Alternatively, \textbf{carefully} transfer the soup in batches to a traditional blender, blend until smooth, and return it to the pot.

      \item \textbf{Finish with Cream \& Cheese} \\
            Place the pot back over low heat. Stir in the heavy whipping cream and Parmesan cheese. Warm gently, stirring continuously, until the cheese is completely melted and the soup is heated through. Do not let it boil, or the cream may curdle. Taste and adjust the seasoning with additional sugar, pepper, or cheese. Serve hot.
\end{enumerate}

\end{document}